\documentclass[11pt]{article}

\usepackage[a4paper,margin=1in]{geometry}
\usepackage[T1]{fontenc}
\usepackage[utf8]{inputenc}
\usepackage{lmodern}
\usepackage{microtype}
\usepackage{amsmath,amssymb,amsthm,mathtools}
\usepackage{enumitem}
\usepackage{xcolor}
\usepackage{graphicx}
\usepackage[hidelinks]{hyperref}
\usepackage[nameinlink,noabbrev]{cleveref}

\setlist{nosep}
\setlength{\parindent}{0pt}
\setlength{\parskip}{0.55em}
\setlength{\emergencystretch}{2em}

\newtheorem{theorem}{Theorem}[section]
\newtheorem{proposition}[theorem]{Proposition}
\newtheorem{lemma}[theorem]{Lemma}
\newtheorem{corollary}[theorem]{Corollary}
\newtheorem{remark}[theorem]{Remark}
\newtheorem{question}[theorem]{Question}
\newtheorem{definition}[theorem]{Definition}

\newcommand{\Cp}{C_p}
\newcommand{\Lp}{L_p}
\newcommand{\wstar}{w^*}
\newcommand{\supp}{\operatorname{supp}}
\newcommand{\spann}{\operatorname{span}}
\newcommand{\cl}{\operatorname{cl}}
\newcommand{\R}{\mathbb{R}}
\newcommand{\N}{\mathbb{N}}
\newcommand{\DeltaK}{\Delta(K)}
\newcommand{\czerop}{(c_0)_p}
\newcommand{\linfyp}{(\ell_\infty)_p}

\title{Weak-Star Basic Sequences in $\Cp(K)'$: \\
Verification, Literature Status, and an Exact Reduction}
\author{Run \texttt{fa\_banach\_001} packet upgrade}
\date{June 22, 2026}

\begin{document}
\maketitle

\begin{abstract}
The supplied argument is correct: if a compact Hausdorff space $K$ contains a
nontrivial convergent sequence, then the finitely supported measures
$\delta_{x_n}-\delta_x$ form a weak-star basic sequence in $\Cp(K)'$.  A
literature search did not locate a proof or counterexample for the full problem
posed as Problem 12 in K\k{a}kol--L\'opez-Pellicer--\'Sliwa (2026).

The main additional result of this report is an exact reduction.  Combining
the known quotient characterization of the finitely supported
Josefson--Nissenzweig property with an elementary norm-rigidity observation, we
prove that, for every infinite compact Hausdorff space $K$, the dual
$\Cp(K)'$ contains a weak-star basic sequence if and only if $\Cp(K)$ has an
infinite-dimensional metrizable quotient.  Thus the 2026 weak-star
basic-sequence question is equivalent to the older metrizable-quotient problem
for $\Cp(K)$.  The same equivalence in fact holds for every infinite
Tychonoff space.  In the compact setting it yields a complete positive answer
for every non-Efimov compactum and for every compactum with the finitely supported
Josefson--Nissenzweig property, but it does not settle the remaining Efimov
case.  The reduction is proved here and was not located in the searched
literature; it should therefore be regarded as an unreviewed deduction.
\end{abstract}

\section*{Packet status}

\textbf{stronger\_partial\_result / exact reduction.}  This packet upgrades
the previous convergent-sequence subcase packet for source arXiv:2601.19873,
Problem 12.  It does not solve the full compact-space problem.  Its promoted
claim is the exact equivalence
\[
\begin{aligned}
  &\Cp(K)' \hbox{ has a weak-star basic sequence}\\
  &\qquad\Longleftrightarrow\qquad
  \Cp(K) \hbox{ has an infinite-dimensional metrizable quotient}.
\end{aligned}
\]
for infinite compact Hausdorff spaces \(K\), with the same equivalence
recorded for arbitrary infinite Tychonoff spaces.  Thus the 2026 question is
reduced to the older metrizable-quotient problem for \(\Cp(K)\).

\medskip

The supplied TeX/PDF report is archived in
\path{evidence/supplied_report_2026_06_22_cp_weakstar_basic/}.  The previous
packet version is preserved in
\path{history/previous_packet_2026_06_22_before_metrizable_quotient_reduction/}.

\section{Problem, notation, and conclusions}

For a Tychonoff space $K$, $\Cp(K)$ denotes $C(K)$ with the topology of
pointwise convergence.  Its continuous dual is
\[
 \Lp(K)=\Cp(K)'=\Delta(K)
 =\left\{\sum_{j=1}^m a_j\delta_{x_j}:m<\infty,\ a_j\in\R,\ x_j\in K\right\},
\]
the vector space of finitely supported signed measures.  It is equipped with
\[
 \sigma(\Lp(K),\Cp(K))=\sigma(\Delta(K),C(K)),
\]
which we call the weak-star topology.

A sequence is weak-star basic if it is a Schauder basis of its weak-star closed
linear span.  In particular, its coefficient functionals must be continuous
for the subspace weak-star topology.

The source paper asks:
\begin{quote}
Does $\Cp(K)'$ admit a weak-star basic sequence for every infinite compact
space $K$?
\end{quote}

\begin{center}
\includegraphics[width=\textwidth]{figures/open_problem_crop.png}
\end{center}

The conclusions reached here are as follows.

\begin{enumerate}[label=\textup{(\roman*)}]
\item The supplied convergent-sequence proof is valid without substantive
      correction; see \cref{sec:verification}.
\item No later proof or counterexample was found in the literature search
      described in \cref{sec:literature}.  The arXiv record of the source still
      has only its January 2026 first version.
\item The problem is positive for every non-Efimov compactum; equivalently, it
      is positive whenever $K$ contains a nontrivial convergent sequence or a
      copy of $\beta\N$.
\item Most importantly, the problem is equivalent to the existence of an
      infinite-dimensional metrizable quotient of $\Cp(K)$; see
      \cref{thm:equivalence}.
\item Both global problems reduce to infinite separable compacta: any infinite
      compactum contains an infinite separable closed subcompactum, and
      restriction transfers a positive answer upward.
\end{enumerate}

\section{Verification of the supplied theorem}\label{sec:verification}

\begin{theorem}[Convergent-sequence case]\label{thm:sequence}
Let $K$ be a compact Hausdorff space.  Suppose that $(x_n)$ is a sequence of
pairwise distinct points converging to
$x\in K\setminus\{x_n:n\in\N\}$.  Then
\[
 \mu_n=\delta_{x_n}-\delta_x\qquad(n\in\N)
\]
is a weak-star basic sequence in $\Lp(K)$.
\end{theorem}

\begin{proof}
Put
\[
 A=\{x\}\cup\{x_n:n\in\N\},\qquad
 M=\spann\{\mu_n:n\in\N\}.
\]
The set $A$ is compact and hence closed in $K$.

We first prove that $M$ is weak-star closed in the \emph{relative} dual
$\Lp(K)$.  Let
\[
 \nu=\sum_{j=1}^m a_j\delta_{t_j}\in\overline M^{\,\wstar},
\]
where the points $t_j$ are distinct and all $a_j$ are nonzero.  If
$t_j\notin A$, normality of $K$ gives $f\in C(K)$ such that
\[
 f(t_j)=1,\qquad
 f\big|_{A\cup(\{t_1,\ldots,t_m\}\setminus\{t_j\})}=0.
\]
Every member of $M$ annihilates $f$, whereas $\nu(f)=a_j\neq0$, a
contradiction.  Hence $\supp\nu\subseteq A$.

Since $\nu$ has finite support, for some $N$ it has the form
\[
 \nu=a\delta_x+\sum_{n=1}^N b_n\delta_{x_n}.
\]
Every member of $M$ annihilates the constant function $1$, and therefore so
does $\nu$.  Thus $a+\sum_{n=1}^N b_n=0$, whence
\[
 \nu=\sum_{n=1}^N b_n(\delta_{x_n}-\delta_x)
     =\sum_{n=1}^N b_n\mu_n\in M.
\]
This proves that $M$ is weak-star closed.

The vectors $\mu_n$ are linearly independent, so every member of $M$ has a
unique finite expansion.  For fixed $n$, the point $x_n$ is isolated in $A$.
Consequently, by normality there exists $f_n\in C(K)$ such that
\[
 f_n(x_n)=1,\qquad f_n\big|_{A\setminus\{x_n\}}=0.
\]
For every finite sum $\sum_k c_k\mu_k$ we then have
\[
 \left(\sum_k c_k\mu_k\right)(f_n)=c_n.
\]
Hence the $n$th coefficient functional is the restriction to $M$ of the
weak-star continuous functional $\nu\mapsto\nu(f_n)$.  Finally, every basis
expansion in $M$ is finite, so its sequence of partial sums is eventually
constant.  Therefore $(\mu_n)$ is a Schauder basis of the weak-star closed
space $M$.
\end{proof}

\begin{remark}[What was checked]
The key point is that closure is taken inside $\Lp(K)$, whose elements are
already finitely supported.  The proof does not claim that $M$ is weak-star
closed in the full Banach dual $C(K)^*$, where limits may be arbitrary Radon
measures.  With this relative-closure convention, every step of the supplied
argument is correct.
\end{remark}

The following elementary formula gives another view of the closed-span
argument.

\begin{lemma}[Closure of point masses]\label{lem:point-closure}
Let $X$ be compact Hausdorff and let $D\subseteq X$.  Inside $\Lp(X)$,
\[
 \overline{\spann\{\delta_d:d\in D\}}^{\,\wstar}
 =\spann\{\delta_x:x\in\cl_X D\}.
\]
\end{lemma}

\begin{proof}
The annihilator of $\spann\{\delta_d:d\in D\}$ in $C(X)$ is the set of
continuous functions vanishing on $D$, hence the set of functions vanishing on
$\cl_XD$.  Taking the annihilator once more gives all finitely supported
measures supported on $\cl_XD$.  Equivalently, one may separate each point
outside $\cl_XD$ from $\cl_XD$ by a continuous function.
\end{proof}

In \cref{thm:sequence}, $M$ is exactly the space of finite measures supported
on $A$ and having total mass zero.  Both conditions are weak-star closed.

\section{Transfer and quotient principles}

\begin{lemma}[Restriction to a closed compact subspace]\label{lem:restriction}
Let $X$ be compact Hausdorff and $K\subseteq X$ closed.  The restriction map
\[
 R:\Cp(X)\longrightarrow\Cp(K),\qquad Rf=f|_K,
\]
is a continuous open surjection.  Its adjoint
\[
 R':\Lp(K)\longrightarrow\Lp(X)
\]
is a weak-star topological embedding with weak-star closed range.  Therefore,
if $\Lp(K)$ has a weak-star basic sequence, then so does $\Lp(X)$.
\end{lemma}

\begin{proof}
Surjectivity follows from the Tietze extension theorem.  To see openness, let
$F\subseteq X$ be finite and let
\[
 U=\{f\in C(X): |f(t)|<\varepsilon_t\text{ for every }t\in F\}
\]
be a basic zero-neighborhood.  Then $R(U)$ contains the basic neighborhood of
zero in $\Cp(K)$ imposing the same inequalities on $F\cap K$.  Indeed, if
$g\in C(K)$ satisfies these inequalities, define a continuous function on
$K\cup(F\setminus K)$ by taking the value $g$ on $K$ and $0$ on
$F\setminus K$, and extend it to $X$.

The assertions about the adjoint are the standard dual consequences of an
open quotient map.  Concretely, its range consists of the finitely supported
measures supported on $K$, which is weak-star closed by point separation.
\end{proof}

\begin{corollary}[Reduction to separable compacta]\label{cor:separable-reduction}
Each of the following global assertions is equivalent to its restriction to
infinite separable compact Hausdorff spaces:
\begin{enumerate}[label=\textup{(\alph*)}]
\item $\Lp(K)$ contains a weak-star basic sequence;
\item $\Cp(K)$ has an infinite-dimensional metrizable quotient.
\end{enumerate}
In particular, if a counterexample to either assertion exists, then a
separable compact counterexample exists.
\end{corollary}

\begin{proof}
Only the upward implication requires proof.  Given an infinite compact $X$,
choose a countably infinite set $D\subseteq X$ and put $K=\cl_XD$.  Then $K$
is infinite, compact, closed in $X$, and separable.  For \textup{(a)}, apply
\cref{lem:restriction}.  For \textup{(b)}, compose the open restriction map
$\Cp(X)\to\Cp(K)$ with an open quotient of $\Cp(K)$.
\end{proof}

\begin{lemma}[Quotients of weak spaces]\label{lem:weak-quotient}
Let $(E,\sigma(E,F))$ be a separated weak locally convex space and let
$N\subseteq E$ be closed.  The quotient topology on $E/N$ is
\[
 \sigma(E/N,N^\perp),
 \qquad
 N^\perp=\{\varphi\in F:\varphi|_N=0\}.
\]
In particular, $E/N$ is metrizable if and only if $N^\perp$ has countable
Hamel dimension.
\end{lemma}

\begin{proof}
Every member of $N^\perp$ descends to a quotient-continuous functional, so the
quotient topology is at least as strong as the displayed weak topology.
Conversely, take a weak neighborhood in $E$ determined by
$\varphi_1,\ldots,\varphi_m\in F$.  Let
$A:E\to\R^m$ be $A(x)=(\varphi_1(x),\ldots,\varphi_m(x))$.
The saturation of that neighborhood by $N$ is the inverse image under $A$ of
an open set saturated by the finite-dimensional subspace $A(N)$.  Passing to
$\R^m/A(N)$ shows that this saturation contains a neighborhood determined by
finitely many linear combinations of the $\varphi_j$ which vanish on $N$.
Thus the quotient topology has a base of $\sigma(E/N,N^\perp)$-neighborhoods.
If $N^\perp$ has a countable Hamel basis, that basis gives a countable
family of seminorms generating the quotient topology.  Conversely, suppose
that the weak quotient is metrizable.  Choose a countable zero-neighborhood
base and, inside each member, a basic weak neighborhood determined by finitely
many elements of $N^\perp$.  The union $G$ of those finite sets is countable.
The family $G$ therefore generates the quotient topology.  For any
$\psi\in N^\perp$, continuity at zero supplies a finite set
$H\subseteq G$ and a constant $C>0$ such that
\[
 |\psi(x+N)|\leq C\max_{\varphi\in H}|\varphi(x)|
 \qquad(x\in E).
\]
Thus $\psi$ vanishes on $\bigcap_{\varphi\in H}\ker\varphi$ and hence, by
finite-dimensional linear algebra, belongs to $\spann H$.  Consequently
$N^\perp=\spann G$.
\end{proof}

\section{Literature status and the known positive region}\label{sec:literature}

The exact source is K\k{a}kol--L\'opez-Pellicer--\'Sliwa,
\emph{On weak*-basic sequences in duals and biduals of spaces $C(X)$ and
Quojections}~\cite{KLS2026}.  Its Problem~12 is the question studied here, and
its Theorem~17 proves that an infinite-dimensional metrizable quotient of
$\Cp(X)$ yields a weak-star basic sequence in $\Cp(X)'$.  The arXiv record
still lists only the version submitted on January~27, 2026.  A June~2, 2026
seminar announcement by K\k{a}kol states both that the two-disjoint-copies
property yields such a quotient and that the metrizable-quotient question for
arbitrary infinite compacta remains a long-standing open problem
\cite{KakolSeminar2026}.

The older paper of Banakh--K\k{a}kol--\'Sliwa proves that if a compact $K$ is
not Efimov, then $\Cp(K)$ has a quotient isomorphic, with the product topology,
to $c_0$ or $\ell_\infty$~\cite[Corollary~3]{BKS2018MQ}.  Since both
$\czerop$ and $\linfyp$ have their canonical Schauder bases in the product
topology, Theorem~17 of~\cite{KLS2026} immediately gives the following.

\begin{corollary}[Non-Efimov compacta]\label{cor:non-efimov}
If an infinite compact Hausdorff space $K$ contains either a nontrivial
convergent sequence or a copy of $\beta\N$, then $\Lp(K)$ has a weak-star
basic sequence.
\end{corollary}

The first alternative is also covered constructively by \cref{thm:sequence}.
In the second alternative, one may use finite block averages on the
$C^*$-embedded copy of $\N$; they are the adjoints of the coordinate
functionals of an open quotient onto $\linfyp$.

A second large positive class comes from the finitely supported
Josefson--Nissenzweig property.

\begin{definition}
A compact space $K$ has the \emph{finitely supported
Josefson--Nissenzweig property} (fsJNP) if there is a sequence
$(\nu_n)\subseteq\Delta(K)$ such that
\[
 \|\nu_n\|=1\quad(n\in\N),
 \qquad
 \nu_n(f)\longrightarrow0\quad(f\in C(K)).
\]
\end{definition}

Banakh--K\k{a}kol--\'Sliwa proved that the fsJNP is equivalent to the
existence of a quotient of $\Cp(K)$ isomorphic to $\czerop$, and also to the
existence of a complemented copy of $\czerop$ in
$\Cp(K)$~\cite[Theorem~1]{BKS2019JNP}.  Consequently:

\begin{corollary}\label{cor:fsjnp}
Every compact space with the fsJNP has a weak-star basic sequence in its
$\Cp$-dual.
\end{corollary}

K\k{a}kol--Sobota--Zdomskyy proved that many classical consistent Efimov
spaces, including inverse limits built from simple extensions, have the
fsJNP~\cite[Corollary~6.12]{KSZ2023}.  They also note that consistently there
are Efimov compacta with the Grothendieck property and hence without the
fsJNP~\cite[p.~16]{KSZ2023}.  Thus neither the Efimov property nor the failure
of the fsJNP is by itself a known obstruction.  Indeed, the June 2026 seminar
announcement reports that the Brech and Sobota--Zdomskyy compacta of Efimov
type have the two-disjoint-copies property and therefore have metrizable
$\Cp$-quotients~\cite{KakolSeminar2026}.

The search covered the exact problem wording, the exact paper title and its
arXiv record, later arXiv papers using the terms ``weak-star basic sequence'',
``$\Cp(X)$'', ``metrizable quotient'', and ``Efimov'', and the citation chains
of the principal 2018--2026 papers.  No full proof or counterexample was found
through June 22, 2026.  This is necessarily a search report rather than a proof
of nonexistence of an unpublished result.

\section{Sequential norm rigidity without the fsJNP}
\label{sec:norm-rigidity}

The absence of the fsJNP has a useful elementary consequence.  It is also the
sequential core of the equivalence between the lack of the fsJNP and the
$\Delta$- and $\ell_1$-Grothendieck properties established in
\cite[Theorem~4.1]{KSZ2023}.

\begin{lemma}[Sequential norm rigidity]\label{lem:norm-rigidity}
Assume that a compact space $K$ does not have the fsJNP.  If
$(\nu_n)\subseteq\Delta(K)$ and $\nu\in\Delta(K)$ satisfy
\[
 \nu_n\longrightarrow\nu
 \quad\hbox{in }\sigma(\Delta(K),C(K)),
\]
then $\|\nu_n-\nu\|\to0$ in total variation norm.
\end{lemma}

\begin{proof}
Otherwise, after passing to a subsequence, there is an $\varepsilon>0$ such that
$\|\nu_n-\nu\|\geq\varepsilon$ for every $n$.  Then
\[
 \lambda_n=\frac{\nu_n-\nu}{\|\nu_n-\nu\|}
\]
is a norm-one sequence in $\Delta(K)$ and $\lambda_n(f)\to0$ for every
$f\in C(K)$, which is precisely an fsJN-sequence.  This contradicts the
assumption.
\end{proof}

The next proposition is the main technical observation of this report.

\begin{proposition}[Basic expansions become norm convergent]\label{prop:norm}
Let $K$ have no fsJNP, and let $(\mu_n)$ be a weak-star basic sequence in
$\Lp(K)$.  Let
\[
 M=\overline{\spann\{\mu_n:n\in\N\}}^{\,\wstar}\subseteq\Lp(K),
 \qquad
 S=\bigcup_{n\in\N}\supp\mu_n.
\]
Then $S$ is countable and every $\mu\in M$ is supported on $S$.  More precisely,
if
\[
 \mu=\sum_{n=1}^{\infty}a_n\mu_n
\]
is its basis expansion, then the partial sums converge to $\mu$ in total
variation norm.  Consequently,
\[
 M\subseteq\Delta(S),
\]
so $M$ has countable Hamel dimension.
\end{proposition}

\begin{proof}
The set $S$ is a countable union of finite sets.  Fix $\mu\in M$ and let
\[
 s_N=\sum_{n=1}^N a_n\mu_n.
\]
By the basis property, $s_N\to\mu$ in
$\sigma(\Delta(K),C(K))$.  Hence \cref{lem:norm-rigidity} gives
$\|s_N-\mu\|\to0$.

Every $s_N$ is supported on $S$.  Since $S$ is countable, it is Borel, and
\[
 \bigl\|\mu|_{K\setminus S}\bigr\|
 \leq \|\mu-s_N\|\longrightarrow0.
\]
Thus $\mu$ is supported on $S$.  Since $\mu\in\Lp(K)=\Delta(K)$ is finitely
supported, $\mu\in\Delta(S)$.  Because $S$ is countable, $\Delta(S)$ has a
countable Hamel basis.
\end{proof}

\begin{corollary}[Metrizable quotient in the no-fsJNP case]
\label{cor:groth-quotient}
If $K$ has no fsJNP and $\Lp(K)$ has a weak-star basic sequence, then
$\Cp(K)$ has an infinite-dimensional metrizable quotient.
\end{corollary}

\begin{proof}
Let $M$ be the weak-star closed span of the basic sequence and put
\[
 N=M_\perp=\{f\in C(K):\mu(f)=0\text{ for all }\mu\in M\}.
\]
By the bipolar theorem for the dual pair $(C(K),\Lp(K))$ and the weak-star
closedness of $M$,
\[
 N^\perp=M.
\]
By \cref{prop:norm}, $M$ has countable Hamel dimension.  Applying
\cref{lem:weak-quotient} to
$\Cp(K)=\bigl(C(K),\sigma(C(K),\Lp(K))\bigr)$, the quotient
$\Cp(K)/N$ is metrizable.  It is infinite-dimensional because its continuous
dual $M$ is infinite-dimensional.
\end{proof}

\section{Exact equivalence with the metrizable-quotient problem}

\begin{theorem}[Main reduction]\label{thm:equivalence}
For every infinite compact Hausdorff space $K$, the following are equivalent.
\begin{enumerate}[label=\textup{(\arabic*)}]
\item $\Lp(K)=\Cp(K)'$ contains a weak-star basic sequence.
\item $\Cp(K)$ has an infinite-dimensional metrizable quotient.
\item $\Cp(K)$ has an infinite-dimensional metrizable quotient possessing a
      Schauder basis.
\item $\Lp(K)$ contains an infinite-dimensional weak-star closed linear
      subspace of countable Hamel dimension.
\end{enumerate}
\end{theorem}

\begin{proof}
First, \textup{(2)} and \textup{(4)} are equivalent.  If
$q:\Cp(K)\to E$ is an infinite-dimensional metrizable quotient and
$N=\ker q$, then \cref{lem:weak-quotient} identifies $E$ with the weak space
$\sigma(C(K)/N,N^\perp)$ and shows that $N^\perp$ has countable Hamel
dimension.  It is infinite-dimensional and weak-star closed in $\Lp(K)$.

Conversely, let $M\subseteq\Lp(K)$ be infinite-dimensional, weak-star closed,
and of countable Hamel dimension.  Put $N=M_\perp$.  Then $N^\perp=M$, and
\cref{lem:weak-quotient} shows that $\Cp(K)/N$ is an infinite-dimensional
metrizable quotient.

The implication \textup{(2)}$\Rightarrow$\textup{(3)} is
Theorem~17 of~\cite{KLS2026}: every infinite-dimensional metrizable quotient
of $\Cp(K)$ has a Schauder basis.  For
\textup{(3)}$\Rightarrow$\textup{(1)}, let $(e_n)$ be such a basis in the
quotient $E$, with coordinate functionals $(e_n^*)$.  If $P_N$ denotes the
$N$th basis projection, then for every $\varphi\in E'$,
\[
 \sum_{n=1}^N \varphi(e_n)e_n^*=\varphi\circ P_N
 \longrightarrow\varphi
 \quad\text{in }\sigma(E',E).
\]
The coefficients are uniquely recovered by evaluation at the $e_n$, so
$(e_n^*)$ is a weak-star Schauder basis of $E'$.  The adjoint of the open
quotient map embeds $E'$ as a weak-star closed subspace of $\Lp(K)$.

Assume \textup{(1)}.  If $K$ has the fsJNP, the theorem of
Banakh--K\k{a}kol--\'Sliwa gives a quotient of $\Cp(K)$ isomorphic to
$\czerop$~\cite[Theorem~1]{BKS2019JNP}, so \textup{(2)} holds.  If $K$
does not have the fsJNP, \cref{cor:groth-quotient} again gives
\textup{(2)}.
\end{proof}

\begin{corollary}[Equivalence of the two global problems]\label{cor:global}
The following assertions are equivalent.
\begin{enumerate}[label=\textup{(\alph*)}]
\item For every infinite compact Hausdorff $K$, $\Cp(K)'$ has a weak-star basic
      sequence.
\item For every infinite compact Hausdorff $K$, $\Cp(K)$ has an
      infinite-dimensional metrizable quotient.
\end{enumerate}
\end{corollary}

Thus Problem 12 of the 2026 paper is not merely related to the earlier
metrizable-quotient problem: it is exactly equivalent to it.

\begin{remark}[Tychonoff extension]\label{rem:tych}
Compactness is not used in the argument of
\cref{lem:norm-rigidity,prop:norm,cor:groth-quotient}.  For an arbitrary
Tychonoff space $X$, the dual of $\Cp(X)$ is still $\Delta(X)$, with
$\|\sum a_j\delta_{x_j}\|=\sum|a_j|$.  Moreover, the JNP--$\czerop$ quotient
characterization~\cite{BKS2019JNP} and Theorem~17 of~\cite{KLS2026} both hold
for Tychonoff spaces.  Hence the four conditions in \cref{thm:equivalence}
are in fact equivalent for every infinite Tychonoff space $X$.
\end{remark}

\section{What a counterexample would have to look like}

Combining \cref{thm:equivalence} with the known positive results gives a sharp
necessary profile.

\begin{corollary}\label{cor:counterexample-profile}
If $K$ is a counterexample to Problem 12, then:
\begin{enumerate}[label=\textup{(\roman*)}]
\item $K$ is an Efimov compactum: it contains neither a nontrivial convergent
      sequence nor a copy of $\beta\N$;
\item $K$ has the $\ell_1$-Grothendieck property, hence no fsJNP;
\item $\Cp(K)$ has no infinite-dimensional metrizable quotient;
\item in particular, $\Cp(K)$ has no quotient isomorphic to $\czerop$ or
      $\linfyp$;
\item $K$ fails the two-disjoint-copies criterion of
      Banakh--K\k{a}kol--\'Sliwa.
\end{enumerate}
Conversely, any infinite compact $K$ satisfying item \textup{(iii)} is a
counterexample.
\end{corollary}

The last converse is simply \cref{thm:equivalence}.  Moreover,
\cref{cor:separable-reduction} shows that, if any counterexample exists, one
can choose a separable one; in fact, the closure of every countably infinite
subset of a counterexample is again a counterexample.  At present, no such
compactum is known from the searched literature.

\section{Why several natural approaches stop short}

\subsection*{A closed copy of $\czerop$ is not enough}
Recent work shows that every infinite compact $K$ supplies natural closed
copies of $\czerop$ in $\Cp(K)$, obtained from pairwise disjoint open sets and
bump functions~\cite{KMS2023}.  But inclusion of a closed subspace
$F\hookrightarrow\Cp(K)$ gives a \emph{surjection}
$\Cp(K)'\to F'$ by restriction, not a canonical weak-star closed embedding of
$F'$ into $\Cp(K)'$.  Such an embedding follows from a continuous projection
onto $F$, and complementability is precisely the extra ingredient supplied by
the fsJNP.  Therefore the closed-subspace theorem alone does not settle the
dual problem.

\subsection*{An infinite metrizable image points in the wrong direction}
Every infinite compact $K$ has an infinite compact metrizable continuous image:
choose countably many continuous functions with infinite-dimensional linear
span and take their diagonal map.  The pullback embeds the corresponding
$\Cp$-space into $\Cp(K)$.  However, a basic sequence in the dual transfers
from a \emph{quotient} or from a closed subspace of the compactum via
\cref{lem:restriction}; an arbitrary continuous image gives the opposite
variance.  A continuous section or averaging operator would repair this, but
such an operator is not available for a general compact surjection.

\subsection*{The product result does not descend}
The complemented-copy theorem for $\Cp(X\times Y)$ gives a weak-star basic
sequence in $\Cp(X\times Y)'$ for infinite $X$ and $Y$~\cite{KMSZ2022}.
Restriction to the diagonal is a quotient
$\Cp(X\times X)\to\Cp(X)$, so its adjoint embeds $\Cp(X)'$ into the product
dual.  This transfers a solution from $X$ to $X\times X$, not from the product
back to $X$.  A basic sequence found somewhere in the larger dual need not lie
in the diagonal-supported range.

\section{Final assessment}

The supplied theorem is correct and gives a particularly transparent explicit
basis.  It was already implicit in the metrizable-quotient machinery and in
earlier $\Cp$-results, so it should be presented as a direct construction rather
than as a solution of the full problem.

No full proof or counterexample was located.  The strongest additional
conclusion obtained here is \cref{thm:equivalence}: weak-star basic sequences
in $\Cp(X)'$ are equivalent to infinite-dimensional metrizable quotients of
$\Cp(X)$, already for arbitrary Tychonoff $X$; see \cref{rem:tych}.  By \cref{cor:separable-reduction}, it is enough to
resolve that quotient problem for separable compacta.  The remaining task is
therefore concentrated entirely in the unresolved separable Efimov case.

\begin{thebibliography}{99}

\bibitem{KLS2026}
J. K\k{a}kol, M. L\'opez-Pellicer, and W. \'Sliwa,
\emph{On weak*-basic sequences in duals and biduals of spaces $C(X)$ and
Quojections}, \href{https://arxiv.org/abs/2601.19873}{arXiv:2601.19873}, 2026.

\bibitem{BKS2018MQ}
T. Banakh, J. K\k{a}kol, and W. \'Sliwa,
\emph{Metrizable quotients of $C_p$-spaces},
Topology Appl. \textbf{249} (2018), 95--102;
\href{https://arxiv.org/abs/1804.02552}{arXiv:1804.02552}.

\bibitem{BKS2019JNP}
T. Banakh, J. K\k{a}kol, and W. \'Sliwa,
\emph{Josefson--Nissenzweig property for $C_p$-spaces},
Rev. R. Acad. Cienc. Exactas F\'is. Nat. Ser. A Mat. RACSAM
\textbf{113} (2019), 3015--3030;
\href{https://arxiv.org/abs/1809.07054}{arXiv:1809.07054}.

\bibitem{KSZ2023}
J. K\k{a}kol, D. Sobota, and L. Zdomskyy,
\emph{Grothendieck $C(K)$-spaces and the Josefson--Nissenzweig theorem},
Fund. Math. \textbf{263} (2023), no.~2, 105--131;
\href{https://arxiv.org/abs/2207.13990}{arXiv:2207.13990}.

\bibitem{KMS2023}
J. K\k{a}kol, A. Molt\'o, and W. \'Sliwa,
\emph{On subspaces of spaces $C_p(X)$ isomorphic to spaces $c_0$ and
$\ell_q$ with the topology induced from $\mathbb{R}^{\mathbb{N}}$},
Rev. Real Acad. Cienc. Exactas Fis. Nat. Ser. A-Mat.
\textbf{117} (2023), article~154.

\bibitem{KMSZ2022}
J. K\k{a}kol, W. Marciszewski, D. Sobota, and L. Zdomskyy,
\emph{On complemented copies of the space $c_0$ in spaces $C_p(X\times Y)$},
Israel J. Math. \textbf{250} (2022), 139--177.

\bibitem{KakolSeminar2026}
J. K\k{a}kol,
\emph{The two-disjoint-copies property for compact spaces, homogeneity, and
relations with the theory of spaces $C_p$ of continuous functions, Part~II},
seminar announcement, Adam Mickiewicz University, June~2, 2026
(\href{https://wmi.amu.edu.pl/zycie-naukowe/wyklady-i-seminaria/seminarium-zakladow-analizy-funkcjonalnej-oraz-analizy-matematycznej/wlasnosc-dwoch-rozlacznych-kopii-dla-przestrzeni-zwartych\%2C-jednorodnosc-i-relacje-z-teoria-przestrzeni-c_p-funkcji-ciaglych-czesc-ii}{online announcement}).

\end{thebibliography}

\end{document}
